\section*{Exercises about Invertibility and Isomorphisms}

\exercise{1}
\begin{xrcs}
  Suppose $T \in \linmap(V,W)$ is invertible. Show that $T^{-1}$ is invertible and $(T^{-1})^{-1} = T$.

  \begin{xprf}
    $T^{-1}$ is the unique element in $\linmap(W,V)$ such that
    \begin{equation}
      T^{-1} T = I_V \quad \text{and} \quad TT^{-1} = I_W.
    \end{equation}
    This equation also means that $T$ satisfies the definition of the inverse of $T^{-1}$.
    Since the inverse of a linear map is unique by Theorem \ref{thm: inverse is unique}, we conclude that $T^{-1}$ is invertible and
    \begin{equation}
      (T^{-1})^{-1} = T.
    \end{equation}
  \end{xprf}
\end{xrcs}

\exercise{2}
\begin{xrcs}
  Suppose $T \in \linmap(U,V)$ and $S \in \linmap(V,W)$ are both invertible linear maps. Prove that $ST \in \linmap(U,W)$ is invertible and that $(ST)^{-1} = T^{-1} S^{-1}$.

  \begin{xprf}
    Let $T \in \linmap(U,V)$ and $S \in \linmap(V,W)$ be invertible linear maps. We compute:
    \begin{equation}
      \begin{aligned}
        (ST) (T^{-1} S^{-1}) &= S (T T^{-1}) S^{-1} = S I_V S^{-1} = S S^{-1} = I_W, \\
        (T^{-1} S^{-1}) (ST) &= T^{-1} (S^{-1} S) T = T^{-1} I_V T = T^{-1} T = I_U.
      \end{aligned}
    \end{equation}
    Since $T^{-1} S^{-1} \in \linmap(W,U)$ is a two-sided inverse of $ST$, we conclude that $ST$ is invertible and $(ST)^{-1} = T^{-1} S^{-1}$.
  \end{xprf}
\end{xrcs}

\exercise{3}
\begin{xrcs}
  Suppose $V$ is finite-dimensional and $T \in \linmap(V)$. Prove that the following are equivalent:
  \begin{enumerate}
    \item $T$ is invertible.
    \item $T v_1, \ddd, T v_n$ is a basis of $V$ for every basis $v_1, \ddd, v_n$ of $V$.
    \item $T v_1, \ddd, T v_n$ is a basis of $V$ for some basis $v_1, \ddd, v_n$ of $V$.
  \end{enumerate}

  \begin{xprf}
    For all proof steps, let $V$ be a finite-dimensional vector space with $\dim V = n$, and let $T \in \linmap(V)$.

    \StepOne Clearly, (b) implies (c).

    \StepTwo Next, we show that (c) implies (a). Suppose $Tv_1, \ddd, Tv_n$ is a basis of $V$ for some basis $v_1, \ddd, v_n$ of $V$.
    Then $\myrange T = \myspan{Tv_1, \dots, Tv_n} = V$, so $T$ is surjective.
    By Theorem \ref{thm: injectivity is equivalent to surjectivity}, a linear operator on a finite-dimensional vector space is surjective if and only if it is invertible.
    Thus, $T$ is invertible.

    \StepThree Now we show that (a) implies (b). Suppose $T$ is invertible, and let $u_1, \dots, u_n$ be any basis of $V$.
    Suppose $c_1 Tu_1 + \dots + c_n Tu_n = 0$ for scalars $c_j \in \myF$.
    Applying the linear map $T^{-1}$ gives
    \begin{equation}
      T^{-1}(c_1 Tu_1 + \dots + c_n Tu_n) = 0 \iff c_1 u_1 + \dots + c_n u_n = 0.
    \end{equation}
    Since $u_1, \dots, u_n$ is a basis, all $c_j = 0$.
    Thus, $Tu_1, \dots, Tu_n$ is a linearly independent list of length $n = \dim V$.
    By Theorem \ref{thm: linearly independent list of the right length is a basis}, $Tu_1, \dots, Tu_n$ is a basis of $V$.
    Hence, (a) implies (b).

    Since (b) $\implies$ (c) $\implies$ (a) $\implies$ (b), all three statements are equivalent.
  \end{xprf}
\end{xrcs}

\exercise{5}
\begin{xrcs}
  Suppose $V$ is finite-dimensional, $U$ is a subspace of $V$, and $S \in \linmap(U,V)$. Prove that there exists an invertible linear map $T \in \linmap(V)$ such that $Tu = Su \quad \forall u \in U \iff S$ is injective.

  \begin{xprf}
    \Rightarrowdirection Suppose there exists an invertible linear map $T \in \linmap(V)$ such that $Tu = Su$ for all $u \in U$.
    If $u_1, u_2 \in U$ satisfy $S(u_1) = S(u_2)$, then $T(u_1) = T(u_2)$.
    Applying $T^{-1}$ gives $u_1 = u_2$. Hence, $S$ is injective.

    \Leftarrowdirection Conversely, suppose $S$ is injective. Let $u_1, \ddd, u_m$ be a basis of $U$.
    Since $S$ is injective, $Su_1, \ddd, Su_m$ is linearly independent in $V$.
    Extend $u_1, \ddd, u_m$ to a basis $u_1, \ddd, u_m, v_{m+1}, \ddd, v_n$ of $V$, and extend $Su_1, \ddd, Su_m$ to a basis $Su_1, \ddd, Su_m, r_{m+1}, \ddd, r_n$ of $V$.
    Define $T \in \linmap(V)$ by specifying its values on the first basis:
    \begin{equation}
      T(u_j) := Su_j \quad (1 \le j \le m) \quad \text{and} \quad T(v_k) := r_k \quad (m+1 \le k \le n).
    \end{equation}
    Then $T(u) = S(u)$ for all $u \in U$.
    Moreover, $T$ maps the basis $u_1, \dots, u_m, v_{m+1}, \dots, v_n$ to the basis $Su_1, \dots, Su_m, r_{m+1}, \dots, r_n$ of $V$.
    By Exercise 3, $T$ is an invertible linear map on $V$.
  \end{xprf}
\end{xrcs}

\exercise{9}
\begin{xrcs}
  Suppose $V$ is finite-dimensional and $S, T \in \linmap(V)$ satisfy $ST = I$. Prove that $T$ is invertible and $T^{-1} = S$ (and hence $TS = I$).
% \begin{xprf}
%   Suppose $v \in V$ satisfies $Tv = 0$. Then
%   \[
%     v = Iv = (ST)v = S(Tv) = S(0) = 0.
%   \]
%   Thus $\mynull T = \{0\}$, which means $T$ is injective.
%   Since $V$ is finite-dimensional, an injective linear operator on $V$ is invertible (Theorem \ref{thm: injectivity is equivalent to surjectivity}).
%   Multiplying the equation $ST = I$ on the right by $T^{-1}$ yields:
%   \[
%     (ST)T^{-1} = I T^{-1} \iff S(TT^{-1}) = T^{-1} \iff S = T^{-1}.
%   \]
%   Therefore, $T$ is invertible with $T^{-1} = S$, and consequently $TS = T T^{-1} = I$.
% \end{xprf}
\end{xrcs}

\exercise{10}
\begin{xrcs}
  Suppose $V$ is finite-dimensional and $T \in \linmap(V)$ commutes with every operator on $V$ (i.e. $TS = ST$ for all $S \in \linmap(V)$). Prove that $T$ is a scalar multiple of the identity operator.
% \begin{xprf}
%   Let $v \in V \setminus \{0\}$. Extend $v$ to a basis $v, v_2, \dots, v_n$ of $V$.
%   Define $S \in \linmap(V)$ by $Sv := v$ and $S v_j := 0$ for each $j \ge 2$.
%   Since $T$ and $S$ commute, $T(Sv) = S(Tv)$.
%   Because $Sv = v$, this gives $Tv = S(Tv)$.
%   The range of $S$ is $\myspan{v}$, so $S(Tv) \in \myspan{v}$, which forces $Tv = \lambda_v v$ for some scalar $\lambda_v \in \myF$.
%   Thus every nonzero vector in $V$ is an eigenvector of $T$.
%   Now let $u, w \in V$ be linearly independent. Then $Tu = \lambda_u u$, $Tw = \lambda_w w$, and $T(u+w) = \lambda_{u+w}(u+w)$.
%   This gives $\lambda_u u + \lambda_w w = \lambda_{u+w} u + \lambda_{u+w} w$, so $(\lambda_u - \lambda_{u+w})u + (\lambda_w - \lambda_{u+w})w = 0$.
%   By linear independence, $\lambda_u = \lambda_{u+w} = \lambda_w$.
%   Hence the eigenvalue $\lambda$ is the same for all vectors in $V$, which proves $Tv = \lambda v$ for all $v \in V$, i.e. $T = \lambda I$.
% \end{xprf}
\end{xrcs}

\exercise{11}
\begin{xrcs}
  Suppose $V$ is finite-dimensional and $S, T \in \linmap(V)$. Prove that $ST$ is invertible if and only if both $S$ and $T$ are invertible.

% \begin{xprf}
%   \Leftarrowdirection If $S$ and $T$ are invertible, then by Exercise 2, $ST$ is invertible and $(ST)^{-1} = T^{-1} S^{-1}$.
%
%   \Rightarrowdirection Conversely, suppose $ST$ is invertible.
%   If $Tv = 0$ for some $v \in V$, then $(ST)v = S(Tv) = S(0) = 0$.
%   Since $ST$ is invertible, $\mynull(ST) = \{0\}$, which forces $v = 0$.
%   Thus $\mynull T = \{0\}$, meaning $T$ is injective.
%   Since $V$ is finite-dimensional, an injective operator on $V$ is invertible (Theorem \ref{thm: injectivity is equivalent to surjectivity}).
%   Now since $T$ is invertible, $T^{-1}$ exists.
%   Multiplying the invertible operator $ST$ on the right by $T^{-1}$ gives $S = (ST) T^{-1}$.
%   Since the product of two invertible operators is invertible, $S$ is invertible.
%   Therefore, both $S$ and $T$ are invertible.
% \end{xprf}
\end{xrcs}
